\documentclass[11pt]{article}
\usepackage[a4paper,margin=1in]{geometry}
\usepackage[T1]{fontenc}
\usepackage[utf8]{inputenc}
\usepackage{lmodern}
\usepackage{amsmath,amssymb,amsthm}
%\usepackage{booktabs}
\usepackage{float}
\usepackage{enumitem}
\usepackage{hyperref}
\usepackage{microtype}
\usepackage{xcolor}
\usepackage{tikz}
\usepackage{standalone}
\usepackage{amsmath}
\usepackage{upgreek}
\usetikzlibrary{calc}


\definecolor{crossedge}{HTML}{B8C0CC}
\definecolor{circulantedge}{HTML}{718096}
\definecolor{matchingedge}{HTML}{E05A47}



\definecolor{circulantedge}{HTML}{718096}
\definecolor{matchingedge}{HTML}{E05A47}

\hypersetup{colorlinks=false}


\newcommand{\floor}[1]{\left\lfloor #1 \right\rfloor}
\newcommand{\ceil}[1]{\left\lceil #1 \right\rceil}
\newcommand{\cl}{\operatorname{cl}}

\numberwithin{equation}{section}
\newtheorem{theorem}{Theorem}[section]
\newtheorem{problem}[theorem]{Problem}
\newtheorem{lemma}[theorem]{Lemma}
\newtheorem{conj}[theorem]{Conjecture}
\newtheorem{proposition}[theorem]{Proposition}
\newtheorem{corollary}[theorem]{Corollary}
%\theoremstyle{definition}
\newtheorem{definition}[theorem]{Definition}
\theoremstyle{plain}
\newtheorem{fact}[theorem]{Fact}
\theoremstyle{definition}
\newtheorem{cons}[theorem]{Construction}
%\theoremstyle{remark}
\newtheorem{remark}[theorem]{Remark}
\theoremstyle{plain}
\newtheorem{claim}{Claim}[section]

\title{Hamiltonicity of 3/2-tough $2K_2$-free graphs}
\author{
Xiaozheng Chen\footnote{School of Mathematics and Statistics, Zhengzhou University, Zhengzhou 450000, P.R. China. Email: cxz@zzu.edu.cn. Supported by the National Natural Science Foundation of China (No. 12301457).}
, Bo Ning\footnote{College of Computer Science, Nankai University, Tianjin 300350, P.R. China. Email: bo.ning@nankai.edu.cn. Partially supported by the National Natural Science Foundation of China (No. 12371350) and the Fundamental Research Funds for the Central Universities, Nankai University (No. 63243151).}}
\date{}


\begin{document}
\maketitle

\begin{abstract}
A graph is $2K_2$-free if it contains no induced subgraph isomorphic to
$2K_2$. Broersma, Patel, and Pyatkin proved in 2014 that every
25-tough $2K_2$-free graph on at least three vertices is Hamiltonian.
Shan reduced the sufficient toughness bound to $3$ in 2020, and
Ota and Sanka reduced it to $2$ in 2022. We prove that
$3/2$-toughness suffices, answering a question of Ota and Sanka.
The bound is best possible.
\end{abstract}

\medskip
\noindent\textbf{Keywords.} Toughness, $2K_2$-free graph, Hamilton cycle, independent set

\medskip
\noindent\textbf{2020 Mathematics Subject Classification.} 05C35, 05C38, 05C75.


\section{Introduction}\label{sec:introduction}

All graphs considered here are finite and simple. A graph is
\emph{Hamiltonian} if it has a cycle containing all its vertices.
For a graph $G$, let $\omega(G)$ denote its number of components.
For $t\geq0$, the graph $G$ is \emph{$t$-tough} if
$|S|\geq t\omega(G-S)$ for every $S\subseteq V(G)$ with
$\omega(G-S)>1$. 
The toughness of $G$, denoted $\uptau(G)$, 
is the maximum value of $t$ for which $G$ is $t$-tough
(for $K_n$ we define $\uptau(G)=\infty$).
In particular,
$G$ is $3/2$-tough precisely when $2|S|\geq3\omega(G-S)$
for every $S\subseteq V(G)$ with $\omega(G-S)>1$.
 It is well known that every Hamiltonian graph is 1-tough. 
In 1973, Chv\'atal \cite{Chvatal1973} proposed the following conjecture.

\begin{conj}[Chv\'atal~\cite{Chvatal1973}]\label{conj:Chvatal}
    There exists a constant $t_0$ such that every $t_0$-tough graph on at least three vertices is Hamiltonian.
\end{conj}

In 2000, Bauer, Broersma and Veldman~\cite{BauerBroersmaVeldman2000} showed
that, for every $0\leq t<9/4$, there exists a non-Hamiltonian
$t$-tough graph. Thus any constant $t_0$ satisfying Chv\'atal's
conjecture must be at least $9/4$.
Toughness conditions for Hamiltonicity have been studied in many
restricted graph classes, including planar graphs, claw-free graphs
and chordal graphs. For background and related results, we refer
the reader to the survey by Bauer, Broersma and
Schmeichel~\cite{BauerBroersmaSchmeichel2006}.

This paper concerns toughness and Hamiltonicity in $2K_2$-free
graphs. A graph is \emph{$2K_2$-free} if it contains no induced
subgraph isomorphic to the disjoint union of two edges.
In 2014, Broersma, Patel, and Pyatkin proved the following theorem.

\begin{theorem}[Broersma, Patel, and Pyatkin~\cite{BroersmaPatelPyatkin2014}]\label{thm:25-tough}
    Every 25-tough $2K_2$-free graph on at least three vertices is Hamiltonian.
\end{theorem}

Shan~\cite{Shan2020} reduced the sufficient
toughness bound to $3$, and Ota and Sanka~\cite{OtaSanka2022}
further reduced it to $2$.

A graph is split if its vertex set can be partitioned into a clique
and an independent set. Every split graph is $2K_2$-free: any two
vertex-disjoint edges have adjacent endpoints in the clique.

Kratsch, Lehel and M\"uller~\cite{KratschLehelMuller1996} proved the following theorem.

\begin{theorem}\label{thm:KratschLehelMuller1996}
    Every $3/2$-tough split graph on at least three vertices is
Hamiltonian.
\end{theorem}


Ota and Sanka~\cite{OtaSanka2022} asked whether $3/2$-toughness   
suffices for $2K_2$-free graphs.
In this paper, we give an affirmative
answer. 
The result extends the Hamiltonicity result of
Kratsch, Lehel and M\"uller~\cite{KratschLehelMuller1996} from
split graphs to the strictly larger class of $2K_2$-free graphs
under the same $3/2$-toughness condition.

\begin{theorem}\label{thm:main}
Every $3/2$-tough $2K_2$-free graph on at least three vertices
is Hamiltonian.
\end{theorem}

The bound is best possible, even for split graphs.
We recall the construction of
Chv\'atal~\cite{Chvatal1973}.
For a positive integer $h$, take pairwise disjoint sets
$S$, $R$,
$T$, with $|S|=h$, $|R|=|T|=2h+1$.
Make every vertex of $S$ adjacent to all other vertices, make $R$
a clique and $T$ independent, and place precisely the matching
edges $r_it_i$ between $R$ and $T$.
The resulting graph $H_h$ is split, with clique $S\cup R$, and
$\uptau(H_h)=3h/(2h+1)=3/2-3/(4h+2)$, which approaches $3/2$
from below as $h\to\infty$.
Moreover, $H_h$ has no $2$-factor. Indeed, a $2$-factor would have $2|T|=4h+2$ edges
incident with $T$, whereas at most $2|S|+|R|=4h+1$ are possible.
Consequently, for every $0\leq t<3/2$, some $H_h$ is a non-Hamiltonian
$t$-tough $2K_2$-free graph.


The organization of this paper is as follows.
In Section~\ref{sec:preliminaries}, we bound the orders of paths and cycles
containing a prescribed independent set. In Section~\ref{sec:maximum-independent}, we give these bounds and Theorem~\ref{thm:KratschLehelMuller1996}; consequently, for every maximum independent set of a non-Hamiltonian counterexample, there exists a cycle of order 
$|V(G)|-1$ containing it. We construct two maximum
independent sets with a common nonempty part $I_0$.
In Section~\ref{sec:proof}, we use path rotations and the corresponding
cycles to give an upper bound on $|N_G(I_0)|$ that contradicts
$3/2$-toughness.

\section{Two lemmas on paths and cycles}\label{sec:preliminaries}

For a graph $G$, we write $|G|=|V(G)|$ and $\alpha(G)$ for its
independence number.
For $X\subseteq V(G)$, let $N_G(X)=\bigcup_{x\in X}N_G(x)$.
We   set $N_X(v)=N_G(v)\cap X$ and $N_X(Y)=N_G(Y)\cap X$,
omitting the subscript $G$ when the graph is fixed.
Two disjoint vertex sets are anticomplete if no edge joins them.
The order of a path or cycle is its number of vertices.
A single vertex is allowed as a path; a cycle has at least three
vertices. A path or cycle $Q$ in $G$ is called \emph{dominating} if
$G-V(Q)$ is edgeless, or equivalently, if
$V(G)\setminus V(Q)$ is independent.
For an oriented path or cycle $P$, we write $x_P^-$ and
$x_P^+$ for the predecessor and successor of $x$, when they exist,
and $X_P^-=\{x_P^-:x\in X,\ x_P^-\text{ exists}\}$, with
$X_P^+$ defined similarly. For a subgraph $H$ of $G$ and
$X\subseteq V(H)$, $N_H(X)$ denotes the union of the neighbourhoods
of vertices of $X$ in $H$, whereas $H[X]$ is the subgraph of $H$
induced by $X$.
Subscripts are omitted when the path or cycle is fixed.
The notation $xPy$ denotes the segment from $x$ to $y$ in the
chosen direction, and $x\overleftarrow{P}y$ denotes the segment
in the reverse direction.

\begin{fact}\label{fact:2K2-property}
   If $G$ is $2K_2$-free, then the following statements hold.
\begin{enumerate}[label=\textup{(\roman*)},nosep]
 \item For every two vertex-disjoint edges \(xy,uv\in E(G)\), at least one of
$xu,\ xv,\ yu,\ yv$ is an edge of \(G\).
    \item For every nonempty independent set $I\subseteq V(G)$ and every two
    nonadjacent vertices $x,y\in V(G)$, the sets $N_I(x)$ and $N_I(y)$
    are comparable by inclusion; that is,
        $N_I(x)\subseteq N_I(y)$ or
        $N_I(y)\subseteq N_I(x)$.
        Consequently, every nonempty independent set
$U\subseteq V(G)\setminus I$ contains a vertex $u$ with
$N_I(u)=N_I(U)$.
\item For every $X\subseteq V(G)$, the graph $G-X$ has at most one
    component containing an edge.
\end{enumerate}
\end{fact}

\begin{proof}
Statement (i) follows immediately from the definition of a
$2K_2$-free graph.

For (ii), suppose to the contrary that $N_I(x)\not\subseteq N_I(y)$ and $N_I(y)\not\subseteq N_I(x)$. Then there exist
$u\in N_I(x)\setminus N_I(y)$ and $v\in N_I(y)\setminus N_I(x)$.
Since $I$ is independent, $uv\notin E(G)$. Moreover,
$xy,uy,xv\notin E(G)$. Hence the four vertices $x,u,y,v$
induce the two disjoint edges $xu$ and $yv$, that is,
an induced $2K_2$, a contradiction.

For (iii), suppose that $G-X$ has two distinct components
$H_1$ and $H_2$ and each contains an edge. Choose
    $x_1y_1\in E(H_1)$ and 
    $x_2y_2\in E(H_2)$.
Since $H_1$ and $H_2$ are distinct components of $G-X$, there is
no edge between $\{x_1,y_1\}$ and $\{x_2,y_2\}$. Thus
$\{x_1,y_1,x_2,y_2\}$ induces a $2K_2$, a contradiction.
\end{proof}

\begin{lemma}\label{lem:predecessor-structure}
Let $G$ be a $2K_2$-free graph and let $J$ be an independent set
in $G$. Consider either an oriented dominating path $P=v_0v_1\cdots v_s$
containing $J$, with $v_0,v_s\notin J$, or an oriented dominating cycle $C$
containing $J$. Let $U$ be the set of vertices outside the chosen
path or cycle. Suppose that $M=N_J(U)\neq\varnothing$,
and choose $u\in U$ with $N_J(u)=M$.
  Set
\[
    D=
    \begin{cases}
        M_P^-\cup\{v_s\},&\text{in the path case},\\
        M_C^-,&\text{in the cycle case},
    \end{cases}
    \qquad
    T=\{x\in D:N_U(x)=\varnothing,\,
                   N_J(x)\setminus M\neq\varnothing\}.
\]
In the path case, assume that no path $P'$ containing $J$ with
both endvertices outside $J$ satisfies either of the following:
\begin{enumerate}[label=\textup{(\alph*)},nosep]
    \item $V(P')=V(P)\cup X$ for some
    $\varnothing\neq X\subseteq U$;
    \item $V(P')=(V(P)\setminus\{x\})\cup\{u\}$ for some $x\in T$.
\end{enumerate}
In the cycle case, impose the corresponding assumptions on cycles
$C'$, with $C,C'$ in place of $P,P'$ and without an endvertex
condition. Then the following statements hold.
\begin{enumerate}[label=\textup{(\roman*)},nosep]
    \item $D\cap J=\varnothing$, $D$ is independent, and $u$ is
    anticomplete to $D$. Every vertex of $U$ has at most one
    neighbour in $D$, and $|D\cap N(U)|\leq1$.
    In the path case, $|D|=|M|+1$ and both endvertices of $P$
    are anticomplete to $U$; in the cycle case, $|D|=|M|$.

    \item For every $x\in T$, we have $N_J(x)\supsetneq M$,
    the predecessor $x^-$ exists, $ux^-\notin E(G)$,
    $x^-\notin M\cup D$, and
    \begin{equation}\label{eq:ND-predecessor}
        N_D(x^-)=\{x\}.
    \end{equation}
\end{enumerate}
\end{lemma}

\begin{proof}
All predecessors and successors are taken on the chosen path or
cycle. Since $J$ is independent, the predecessors of vertices of
$M$ lie outside $J$. The predecessor map is injective; in the path
case, $v_s$ lies outside $J$ and is not a predecessor. This proves
$D\cap J=\varnothing$ and the assertions about $|D|$.

In the path case, neither endvertex has a neighbour in $U$, since
such a neighbour would extend $P$, contrary to (a). If $x\in D$
has a successor and $ux\in E(G)$, inserting $u$ into the edge
$xx^+$    contradicts (a), because $x^+\in M=N_J(u)$.
Thus $u$ is anticomplete to $D$ in both cases.

Suppose that distinct vertices $x,z\in D$ are adjacent. In the
path case, assume that $x$ precedes $z$. If $z\neq v_s$, the path
$v_0P xz\overleftarrow P x^+uz^+P v_s$ has vertex set
$V(P)\cup\{u\}$; if $z=v_s$, use
$v_0P xv_s\overleftarrow P x^+u$ instead. In the cycle case,
$xz\overleftarrow C x^+uz^+C x$ has vertex set $V(C)\cup\{u\}$.
In each construction the indicated old segments are disjoint and
partition the original vertex set. Each therefore contradicts (a),
so $D$ is independent.

Suppose that $y\in U$ has two neighbours $x,z\in D$. Then $y\neq u$.
In the path case, neither neighbour is an endvertex; ordering $x$
before $z$, the path $v_0P xyz\overleftarrow P x^+uz^+P v_s$
has vertex set $V(P)\cup\{u,y\}$. In the cycle case, use
$xyz\overleftarrow C x^+uz^+C x$, whose vertex set is
$V(C)\cup\{u,y\}$. Both contradict (a). Hence every vertex of
$U$ has at most one neighbour in $D$. Since $U$ and $D$ are
independent, Fact~\ref{fact:2K2-property}\textup{(ii)} implies that
the sets $N_D(y)$, $y\in U$, are linearly ordered by inclusion.
It follows that $|D\cap N(U)|\leq1$, proving (i).

Let $x\in T$. Since $xu\notin E(G)$,
Fact~\ref{fact:2K2-property}\textup{(ii)} and the definition of $T$
give $N_J(x)\supsetneq N_J(u)=M$. In the path case, $x\neq v_0$:
otherwise $uv_1P v_s$ would be a path excluded by (b).
Thus $x^-$ exists in either case.

If $ux^-\in E(G)$, then in the path case use
$v_0P x^-ux^+P v_s$ when $x\neq v_s$, and $v_0P v_s^-u$
when $x=v_s$. In the cycle case, replace $x^-xx^+$ by $x^-ux^+$.
Each operation replaces $x$ by $u$, contrary to (b).
Consequently $ux^-\notin E(G)$ and $x^-\notin M$.
Also $x^-\notin D$, since its successor is $x\notin J$, whereas
every vertex of $D$ other than $v_s$ in the path case, and every
vertex of $D$ in the cycle case, has its successor in $M$.

Suppose that $x^-z\in E(G)$ for some $z\in D\setminus\{x\}$.
In the path case, one of the following paths replaces $x$ by $u$:
\[
\begin{cases}
 v_0P x^-z\overleftarrow P x^+uz^+P v_s,
   &x,z\neq v_s\text{ and }x\text{ precedes }z,\\
 v_0P zx^-\overleftarrow P z^+ux^+P v_s,
   &x,z\neq v_s\text{ and }z\text{ precedes }x,\\
 v_0P x^-v_s\overleftarrow P x^+u,
   &z=v_s,\\
 v_0P zv_s^-\overleftarrow P z^+u,
   &x=v_s.
\end{cases}
\]
The relation $x^-\notin M\cup D$ ensures that the indicated old
segments are disjoint and partition $V(P)\setminus\{x\}$.
Each path contains $J$ and has both endvertices outside $J$.
In the cycle case, the corresponding cycle is
$x^-z\overleftarrow C x^+uz^+C x^-$: its old arcs $x^+Cz$ and
$z^+Cx^-$ partition $V(C)\setminus\{x\}$. All these constructions
contradict (b). Since $x^-x\in E(G)$, this proves
\eqref{eq:ND-predecessor} and hence (ii).
\end{proof}

\begin{lemma}\label{lem:path}
Let $G$ be a $2K_2$-free graph and let $J$ be an independent set
in $G$. Suppose that $V(G)$ is covered by a family $\mathcal P$
of pairwise vertex-disjoint paths, each with both endvertices
outside $J$. A single vertex outside $J$ is allowed as a path.
Then the following statements hold.
\begin{enumerate}[label=\textup{(\roman*)},nosep]
    \item If $G$ is nonempty, then it has a path containing
    $J$ with both endvertices outside $J$.

    \item If $G$ is nonempty and $\ell$ is the maximum order
    of such a path, then
    \begin{equation}\label{eq:path-bound}
        \alpha(G)\geq |J|+|G|-\ell.
    \end{equation}

    \item The vertex set $V(G)$ can be covered by at most
    $\alpha(G)-|J|+1$ pairwise vertex-disjoint paths, each with
    both endvertices outside $J$.
\end{enumerate}
\end{lemma}

\begin{proof}
If $G$ is empty, (iii) follows by taking the empty family, while
(i) and (ii) do not apply. Hence assume that $G$ is nonempty.

\noindent
(i) The assertion is immediate if $J=\varnothing$.
Otherwise, discard the paths in $\mathcal P$ disjoint from $J$,
orient each remaining path, and shorten it to the segment from the
predecessor of its first vertex in $J$ to the successor of its last
vertex in $J$. Let $\mathcal P_J$ be the resulting family.
Its paths are vertex-disjoint, their intersections with $J$
partition $J$, and each endvertex lies outside $J$ and is
adjacent along its path to a vertex of $J$.

If $|\mathcal P_J|\geq2$, choose distinct paths
$P_1=v_0v_1\cdots v_s$ and $P_2=u_0u_1\cdots u_t$ in this
family. Then $v_1,u_1\in J$ and $v_0,v_s,u_0,u_t\notin J$.
The edges $v_0v_1,u_0u_1$ are disjoint, and $v_1u_1\notin E(G)$.
By $2K_2$-freeness, at least one of $v_0u_0,v_0u_1,u_0v_1$
is an edge. Accordingly, let
\[
P_3=
\begin{cases}
 v_s\overleftarrow {P_1}v_0u_0P_2u_t,
   &v_0u_0\in E(G),\\
 v_s\overleftarrow {P_1}v_0u_1P_2u_t,
   &v_0u_0\notin E(G),\ v_0u_1\in E(G),\\
 v_s\overleftarrow {P_1}v_1u_0P_2u_t,
   &\text{otherwise}.
\end{cases}
\]
The set $(V(P_1)\cup V(P_2))\setminus V(P_3)$ is either empty,
$\{v_0\}$, or $\{u_0\}$; in every case it is disjoint from $J$. The new endvertices are $v_s,u_t$,
and their incident path edges are unchanged. Replacing $P_1,P_2$
by $P_3$ therefore preserves all the stated properties and reduces
the number of paths by one. Repeating this operation gives a single
path containing $J$, proving (i).

\medskip
\noindent
(ii) Let $\ell$ be the maximum order of a path containing $J$
with both endvertices outside $J$.

\begin{claim}\label{claim:U-independent}
Every such path of order $\ell$ is dominating.
\end{claim}

\begin{proof}
Let $P=v_0v_1\cdots v_{\ell-1}$ be such a path and   set
$U=V(G)\setminus V(P)$. If $\ell=1$, then $J=\varnothing$
and $G$ is edgeless. Otherwise, suppose that $xy\in E(G[U])$.
If $v_0$ has a neighbour among $x,y$, say $x$, then
$yxv_0P v_{\ell-1}$ is a longer path of the required type.
If not, $2K_2$-freeness applied to $v_0v_1,xy$ gives, after
interchanging $x,y$ if necessary, $v_1x\in E(G)$. The path
$yxv_1P v_{\ell-1}$ has order $\ell+1$ and still contains $J$
with both endvertices outside $J$, since $v_0\notin J$.
Both cases contradict the choice of $\ell$.
\end{proof}

If $J=\varnothing$, take a longest path $P$ and an endvertex $v_0$.
The vertex $v_0$ has no neighbour outside $P$, so the claim gives
$\alpha(G)\geq |G|-\ell+1$. Hence \eqref{eq:path-bound} holds.

Suppose now that $J\neq\varnothing$ and, to the contrary, that
$\alpha(G)<|J|+|G|-\ell$. Then $\ell<|G|$.
Let $\mathcal P_0$ consist of the dominating paths $P$ containing $J$
with both endvertices outside $J$.
By Claim~\ref{claim:U-independent}, this family
contains a path of order $\ell$.
Choose $P_0\in\mathcal P_0$ first with maximum order and then
with $|N_J(V(G)\setminus V(P_0))|$ maximum.   Set
$U_0=V(G)\setminus V(P_0)$ and $m_0=|N_J(U_0)|$.
For $r\geq1$, once $m_0,\ldots,m_{r-1}$ are defined, let
$\mathcal P_r$ consist of the paths $P\in\mathcal P_0$ satisfying
\begin{equation}\label{eq:path-count}
    \bigl|\{v\in V(G)\setminus V(P):|N_J(v)|>m_i\}\bigr|
    \geq r+1-i \qquad (0\leq i<r).
\end{equation}
Whenever $\mathcal P_r\neq\varnothing$, choose $P_r$ first with
maximum order in $\mathcal P_r$ and then with
$|N_J(V(G)\setminus V(P_r))|$ maximum.   Set
$U_r=V(G)\setminus V(P_r)$ and $m_r=|N_J(U_r)|$.
We prove inductively that
\begin{equation}\label{eq:path-order}
    \mathcal P_r\neq\varnothing,\qquad
    |P_r|=\ell-r,\qquad m_0<m_1<\cdots<m_r.
\end{equation}
The assertion for $r=0$ follows from the initial choice.

Assume that \eqref{eq:path-order} holds for all indices up to $r$.
Write $P=P_r=v_0v_1\cdots v_s$, $U=U_r$, and $M=N_J(U)$.
If $M=\varnothing$, then $J\cup U$ is independent of order
$|J|+|G|-\ell+r$, contrary to our assumption.
Thus $M\neq\varnothing$. By
Fact~\ref{fact:2K2-property}\textup{(ii)}, choose $u\in U$ with
$N_J(u)=M$.   Set $D=M_P^-\cup\{v_s\}$ and
$T=\{x\in D:N_U(x)=\varnothing,\,
N_J(x)\setminus M\neq\varnothing\}$.
In particular, $|N_J(u)|=|M|=m_r$.

We verify the two assumptions of
Lemma~\ref{lem:predecessor-structure}.
First, suppose that a path $P'$ containing $J$ with both
endvertices outside $J$ has vertex set $V(P)\cup X$, where
$\varnothing\neq X\subseteq U$, and   set $k=|X|$.
If $k>r$, then $|P'|=\ell-r+k>\ell$, a contradiction.
Otherwise   set $t=r-k$. For every $i<t$, at least
$r+1-i-k=t+1-i$ vertices outside $P'$ have more than $m_i$
neighbours in $J$. Hence $P'\in\mathcal P_t$ and
$|P'|=\ell-t=|P_t|$.
Moreover, \eqref{eq:path-count} for $P_r$ with $i=t$ gives at
least $k+1$ vertices of $U$ having more than $m_t$ neighbours
in $J$. At most $k$ lie in $X$, so one remains outside $P'$,
contrary to the choice of $m_t$.

Second, for $x\in T$, the nonedge $xu$ and
Fact~\ref{fact:2K2-property}\textup{(ii)} give
$N_J(x)\supsetneq M$. A path $P'$ of the required type with
vertex set $(V(P)\setminus\{x\})\cup\{u\}$ would be dominating,
since $V(G)\setminus V(P')=(U\setminus\{u\})\cup\{x\}$
and $N_U(x)=\varnothing$.
Replacing $u$ by $x$ in the set of vertices outside the path leaves every count in
\eqref{eq:path-count} unchanged, since
$|N_J(x)|>m_r=|N_J(u)|>m_i$ for $i<r$.
Thus $P'\in\mathcal P_r$ and $|P'|=|P_r|$, but the vertex $x$,
which lies outside $P'$, has more than $m_r$ neighbours in $J$,   a
contradiction. Lemma~\ref{lem:predecessor-structure} therefore
applies to $P$.

\medskip
\noindent
\textbf{Case 1: $|T|\leq1$.}

We show that at most one vertex of $D$ belongs to $N(U)\cup T$.
This follows from Lemma~\ref{lem:predecessor-structure}(i) when
$T=\varnothing$. Suppose that $T=\{x\}$ and that $z\in D$
has a neighbour $y\in U$. Then $z\neq x,v_s$ and $y\neq u$.
The edges $x^-x,zy$ have no cross-edges $x^-z,xz,xy$, by
\eqref{eq:ND-predecessor}, the independence of $D$, and the
definition of $T$. Thus $x^-y\in E(G)$ by $2K_2$-freeness.
Define
\[
P'=
\begin{cases}
 v_0P x^-yz\overleftarrow P x^+uz^+P v_s,
   &x\text{ precedes }z,\\
 v_0P zyx^-\overleftarrow P z^+ux^+P v_s,
   &z\text{ precedes }x\neq v_s,\\
 v_0P zyv_s^-\overleftarrow P z^+u,
   &x=v_s.
\end{cases}
\]
The indicated old segments partition $V(P)\setminus\{x\}$.
Thus $P'$ contains $J$, has both endvertices outside $J$, and
has vertex set $(V(P)\setminus\{x\})\cup\{u,y\}$.
Moreover, $P'$ is dominating, since
$V(G)\setminus V(P')=(U\setminus\{u,y\})\cup\{x\}$
and $N_U(x)=\varnothing$.

If $r=0$, then $|P'|=\ell+1$, a contradiction. If $r\geq1$,
then for every $i<r-1$, passing from $U$ to
$(U\setminus\{u,y\})\cup\{x\}$ decreases the count in \eqref{eq:path-count} by
at most one. At least $r-i=(r-1)+1-i$ such vertices remain,
so $P'\in\mathcal P_{r-1}$ and $|P'|=|P_{r-1}|$.
However, $x$ lies outside $P'$ and satisfies
$|N_J(x)|>m_r>m_{r-1}$, contrary to the choice of $m_{r-1}$.
Consequently $D\cap N(U)=\varnothing$ when $T=\{x\}$,
and in either case at most one vertex of $D$ belongs to
$N(U)\cup T$.

The set
$U\cup\bigl(D\setminus(N(U)\cup T)\bigr)\cup(J\setminus M)$
is independent: the middle part is anticomplete to both $U$
and $J\setminus M$, and $N_J(U)=M$.
Its order is at least
$|U|+(|D|-1)+|J\setminus M|=|U|+|J|$.
Since $|U|=|G|-\ell+r$, this contradicts our assumption.

\medskip
\noindent
\textbf{Case 2: $|T|\geq2$.}

Choose distinct $x,z\in T$, with $x$ preceding $z$ on $P$.
By Lemma~\ref{lem:predecessor-structure}, their predecessors
lie outside $M\cup D$, and the edges $x^-x,z^-z$ have no
cross-edges $xz,x^-z,z^-x$.
Thus $x^-z^-\in E(G)$ by $2K_2$-freeness. Let
\[
P'=
\begin{cases}
 v_0P x^-z^-\overleftarrow P x^+uz^+P v_s,
   &z\neq v_s,\\
 v_0P x^-v_s^-\overleftarrow P x^+u,
   &z=v_s.
\end{cases}
\]
The old segments partition $V(P)\setminus\{x,z\}$.
Hence $P'$ contains $J$, has both endvertices outside $J$,
and has order $|P_r|-1$. It is dominating, since
$V(G)\setminus V(P')=(U\setminus\{u\})\cup\{x,z\}$,
$U,D$ are independent, and $x,z$ are anticomplete to $U$.

For every $i<r$, passing from $U$ to $(U\setminus\{u\})\cup\{x,z\}$
increases the count in \eqref{eq:path-count} by one, giving
at least $r+2-i$ vertices above the threshold $m_i$.
For $i=r$, the vertices $x,z$ provide the two required vertices,
since $|N_J(x)|,|N_J(z)|>m_r$. Thus
$P'\in\mathcal P_{r+1}$, so $\mathcal P_{r+1}\neq\varnothing$.
As $\mathcal P_{r+1}\subseteq\mathcal P_r$, its maximum path
order is at most $|P_r|$. Equality would contradict the choice
of $m_r$, since every $Q\in\mathcal P_{r+1}$ has a vertex
outside $Q$ with more than $m_r$ neighbours in $J$.
The path $P'$ attains order $|P_r|-1$, so
$|P_{r+1}|=\ell-(r+1)$. The new threshold    gives
$m_{r+1}>m_r$, completing the induction.

Under our contrary assumption, Case~1 is impossible, so Case~2
proves the induction step for every $r$. This produces an infinite strictly
increasing sequence $m_0<m_1<\cdots$, although $m_r\leq|J|$
for every $r$. This contradiction proves \eqref{eq:path-bound}.

\medskip
\noindent
(iii) Take a path of order $\ell$ as in (ii), together with the
trivial paths on all vertices outside this path. These form a spanning family
of $|G|-\ell+1$ vertex-disjoint paths with both endvertices outside
$J$. By \eqref{eq:path-bound}, their number is at most
$\alpha(G)-|J|+1$.
\end{proof}

\begin{lemma}\label{lem:cycle}
Let $G$ be a $2K_2$-free graph and let $J$ be an independent set
in $G$. Suppose that $G$ has a dominating cycle containing $J$.
If $\ell$ is the maximum order of such a cycle, then
\begin{equation}\label{eq:cycle-bound}
    \alpha(G)\geq |J|+|G|-\ell-1.
\end{equation}
\end{lemma}

\begin{proof}
Let $\mathcal F_0$ be the family of dominating cycles containing $J$.
Choose $C_0\in\mathcal F_0$ first with maximum order and then
with $|N_J(V(G)\setminus V(C_0))|$ maximum.   Set
$\ell_0=|C_0|=\ell$, $U_0=V(G)\setminus V(C_0)$, and
$m_0=|N_J(U_0)|$.
For $r\geq1$, once $m_0,\ldots,m_{r-1}$ are defined, let
$\mathcal F_r$ consist of the cycles $C\in\mathcal F_0$ satisfying
\begin{equation}\label{eq:cycle-count}
    \bigl|\{v\in V(G)\setminus V(C):|N_J(v)|>m_i\}\bigr|
    \geq r+1-i \qquad (0\leq i<r).
\end{equation}
Whenever $\mathcal F_r\neq\varnothing$, choose $C_r$ first
with maximum order in $\mathcal F_r$ and then with
$|N_J(V(G)\setminus V(C_r))|$ maximum.   Set
$\ell_r=|C_r|$, $U_r=V(G)\setminus V(C_r)$, and $m_r=|N_J(U_r)|$.
We argue by induction: either the desired bound follows from an
independent set, or the next extremal cycle satisfies
\begin{equation}\label{eq:cycle-order}
    \ell_r=\ell_0-r,\qquad
    m_0<m_1<\cdots<m_r\leq|J|.
\end{equation}

Assume that \eqref{eq:cycle-order} holds for all indices up to $r$,
and write $C=C_r$, $U=U_r$, and $M=N_J(U)$.
If $M=\varnothing$, then $J\cup U$ is independent and gives the
required bound, since $|U|=|G|-\ell_0+r$. Otherwise, choose $u\in U$ with $N_J(u)=M$ by
Fact~\ref{fact:2K2-property}\textup{(ii)}. Orient $C$ and   set
$A=M_C^-$ and
$T=\{x\in A:N_U(x)=\varnothing,\,
N_J(x)\setminus M\neq\varnothing\}$.

%We first verify the two extremal conditions needed for
%Lemma~\ref{lem:predecessor-structure}.

Suppose that a cycle $C'$ satisfies
    $V(C')=V(C)\cup X$
for some nonempty $X\subseteq U$, and   set $k=|X|$.
If $k>r$, then
    $|C'|=\ell_0-r+k>\ell_0$,
contrary to the choice of $C_0$.
Thus $1\leq k\leq r$.   Set $t=r-k$.
For every $i<t$, the set $V(G)\setminus V(C')=U\setminus X$
contains at least
   $ r+1-i-k=t+1-i$
vertices with more than $m_i$ neighbours in $J$.
Hence $C'\in\mathcal F_t$, and
    $|C'|=\ell_0-t=\ell_t$.
Moreover, applying \eqref{eq:cycle-count} to $C_r$ with $i=t$
shows that at least $k+1$ vertices of $U$ have more than $m_t$
neighbours in $J$. Since at most $k$ of them belong to $X$,
at least one remains outside $C'$. Thus
    $|N_J(V(G)\setminus V(C'))|>m_t$,
contrary to the choice of $C_t$. Therefore no such cycle $C'$
exists.

Now let $x\in T$. Since $xu\notin E(G)$ and
$N_J(x)\setminus M\neq\varnothing$,
Fact~\ref{fact:2K2-property}\textup{(ii)} gives
    $N_J(x)\supsetneq M$,
and hence $|N_J(x)|>m_r$.
Suppose that a cycle $C'$ satisfies
   $ V(C')=(V(C)\setminus\{x\})\cup\{u\}$.
The cycle $C'$ is dominating, since
$V(G)\setminus V(C')=(U\setminus\{u\})\cup\{x\}$
and $N_U(x)=\varnothing$.
For every $i<r$, both $u$ and $x$ have more than $m_i$
neighbours in $J$, so replacing $u$ by $x$ does not decrease
any of the counts in \eqref{eq:cycle-count}.
Thus $C'\in\mathcal F_r$ and $|C'|=\ell_r$.
Since $x\notin V(C')$ and $N_J(x)\supsetneq M$, however,
    $|N_J(V(G)\setminus V(C'))|>m_r$,
contrary to the choice of $C_r$.

The two assumptions of
Lemma~\ref{lem:predecessor-structure} are therefore satisfied.
Applying its cycle case with $D=A$ gives
$|A|=|M|$ together with all the stated independence and
predecessor properties of $A$ and $T$.

\medskip
\noindent
\textbf{Case 1: $|T|\leq1$.}

If $T=\varnothing$, then $|A\cap N(U)|\leq1$.
Suppose that $T=\{x\}$ and that $z\in A$ has a neighbour
$y\in U$. Then $z\neq x$ and $y\neq u$.
By Lemma~\ref{lem:predecessor-structure}, the edges $x^-x,zy$
have no cross-edges $x^-z,xz,xy$, so $x^-y\in E(G)$.
The cycle
 $   C'=x^-yz\overleftarrow C x^+uz^+C x^-$
has vertex set $(V(C)\setminus\{x\})\cup\{u,y\}$:
its old arcs $x^+Cz$ and $z^+Cx^-$ partition $V(C)\setminus\{x\}$.
Moreover, $C'$ is dominating, since
$V(G)\setminus V(C')=(U\setminus\{u,y\})\cup\{x\}$
and $N_U(x)=\varnothing$.
If $r=0$, then $|C'|=\ell_0+1$, a contradiction.
Suppose that $r\ge1$. For each $i<r-1$, both $u$ and $x$ have
more than $m_i$ neighbours in $J$. Thus replacing $U$ by
$(U\setminus\{u,y\})\cup\{x\}$ decreases the count in \eqref{eq:cycle-count}
by at most one. At least $r-i=(r-1)+1-i$ vertices remain above
the threshold $m_i$, so $C'\in\mathcal F_{r-1}$.
Moreover, $|C'|=\ell_r+1=\ell_{r-1}$, whereas
 $|N_J(V(G)\setminus V(C'))|
 \ge |N_J(x)|>m_r>m_{r-1}$.
This contradicts the second extremal choice of $C_{r-1}$.

Hence at most one vertex of $A$ belongs to $N(U)\cup T$.
As in the path case, the set
$U\cup\bigl(A\setminus(N(U)\cup T)\bigr)\cup(J\setminus M)$
is independent. Its order is at least
$|U|+(|A|-1)+|J\setminus M|=|J|+|U|-1$.
This gives the required bound because $|U|=|G|-\ell_0+r$.

\medskip
\noindent
\textbf{Case 2: $|T|\geq2$.}

Choose distinct $x,z\in T$.
Lemma~\ref{lem:predecessor-structure} and $2K_2$-freeness give
$x^-z^-\in E(G)$, exactly as in the path case. The six vertices
$x^-,x,x^+,z^-,z,z^+$ are distinct and occur in this cyclic order.
Consequently
    $C'=x^-z^-\overleftarrow C x^+uz^+C x^-$
is a cycle with vertex set $(V(C)\setminus\{x,z\})\cup\{u\}$:
its old arcs $x^+Cz^-$ and $z^+Cx^-$ partition
$V(C)\setminus\{x,z\}$. The cycle $C'$ is dominating, since
$V(G)\setminus V(C')=(U\setminus\{u\})\cup\{x,z\}$,
$U,A$ are independent, and $x,z$ are anticomplete to $U$.

For each $i<r$, the vertex $u$ is replaced outside the cycle by
$x,z$, and all three vertices have more than $m_i$ neighbours in $J$.
Thus the count in \eqref{eq:cycle-count} increases by one and is at least $r+2-i$.
For the new threshold $m_r$, both $x$ and $z$ have more than $m_r$
neighbours in $J$, giving the two required vertices.
Consequently $C'\in\mathcal F_{r+1}$, so this family is nonempty.
Since $\mathcal F_{r+1}\subseteq\mathcal F_r$, its maximum cycle
order is at most $\ell_r$. If equality held, a maximum-order cycle
$Q\in\mathcal F_{r+1}$ would also have maximum order in
$\mathcal F_r$, but the defining condition for $\mathcal F_{r+1}$
at $i=r$ would imply
$ |N_J(V(G)\setminus V(Q))|>m_r$,
contrary to the choice of $C_r$.
Together with $|C'|=\ell_r-1$, this gives
$\ell_{r+1}=\ell_r-1$. The same condition, applied to $C_{r+1}$,
gives $m_{r+1}>m_r$, completing the induction.

Case~2 cannot apply for every $r$, since it would give an infinite
strictly increasing sequence of integers $m_r\leq|J|$. Hence, for
some $r$, either $M=\varnothing$ or Case~1 gives
\[
    \alpha(G)\geq |J|+|U_r|-1
    =|J|+|G|-\ell_0+r-1
    \geq |J|+|G|-\ell-1,
\]
which proves \eqref{eq:cycle-bound}.
\end{proof}

\begin{corollary}\label{claim:cycle-restricted}
Let $G$ be a $2K_2$-free graph, let $J$ be an independent set,
and let $\mathcal H$ be a nonempty family of dominating cycles
containing $J$.

Starting with $\mathcal F_0=\mathcal H$, define the families
$\mathcal F_r$ and the numbers $m_r$ successively.
For each nonempty $\mathcal F_r$, choose $C_r\in\mathcal F_r$
first with maximum order and then with
$|N_J(V(G)\setminus V(C_r))|$ maximum, and set
 $U_r=V(G)\setminus V(C_r), m_r=|N_J(U_r)|$.
For $r\ge1$, once $m_0,\ldots,m_{r-1}$ are defined,
let $\mathcal F_r$ consist of the cycles $C\in\mathcal H$ satisfying
\[
 \bigl|\{v\in V(G)\setminus V(C):|N_J(v)|>m_i\}\bigr|
 \ge r+1-i\qquad(0\le i<r).
\]

For each selected $C_r$ with $M=N_J(U_r)\ne\emptyset$,
write $C=C_r$ and $U=U_r$, choose $u\in U$ with $N_J(u)=M$,
orient $C$, and set
\[
 A=M_C^-,\qquad
 T=\{x\in A:N_U(x)=\emptyset,\;N_J(x)\setminus M\ne\emptyset\}.
\]
Suppose that at each such stage the following cycles of $G$,
whenever they exist, belong to $\mathcal H$:
\begin{enumerate}[nosep]
\item[(a)] Every cycle with vertex set $V(C)\cup X$, where
$\emptyset\ne X\subseteq U$.

\item[(b)] Every cycle with vertex set
$(V(C)\setminus\{x\})\cup\{u\}$, where $x\in T$.

\item[(c)] The cycle
 $x^-yz\overleftarrow{C}x^+uz^+Cx^-$,
when $T=\{x\}$ and vertices $z\in A\setminus\{x\}$ and
$y\in U\setminus\{u\}$ satisfy $zy\in E(G)$.

\item[(d)] The cycle
$ x^-z^-\overleftarrow{C}x^+uz^+Cx^-$,
for distinct $x,z\in T$.
\end{enumerate}
All predecessors, successors and segments in (c) and (d) are taken
on $C$. If $\ell_{\mathcal H}$ is the maximum order of a cycle in
$\mathcal H$, then
 $\alpha(G)\ge |J|+|G|-\ell_{\mathcal H}-1$.
\end{corollary}

\begin{proof}
Repeat the induction in Lemma~\ref{lem:cycle} with
$\mathcal F_0=\mathcal H$ and $\ell_0=\ell_{\mathcal H}$.
At each stage with $M\ne\emptyset$, conditions (a) and (b)
ensure that the first two extremal comparisons take place in
$\mathcal H$. The order comparisons and the threshold counts in \eqref{eq:cycle-count}
exclude these cycles exactly as in that proof.
These comparisons establish the two hypotheses of Lemma~\ref{lem:predecessor-structure},
which can therefore be applied with $D=A$.

With the independence and predecessor properties from Lemma~\ref{lem:predecessor-structure}
established, conditions (c) and (d) ensure that the cycles constructed
in Cases~1 and~2 also belong to $\mathcal H$.
The same argument consequently gives either an independent set
of order at least $|J|+|G|-\ell_{\mathcal H}-1$, or a nonempty
family $\mathcal F_{r+1}$ with
\[
 |C_{r+1}|=|C_r|-1,\qquad m_{r+1}>m_r.
\]
The latter alternative cannot persist, since $m_r\le |J|$.
The case $M=\emptyset$ is handled by the independent set $J\cup U$,
as in Lemma~\ref{lem:cycle}. This proves the inequality.
\end{proof}

In particular, suppose that every vertex outside a cycle must lie
in a fixed set. A cycle with vertex set $V(C_r)\cup X$ preserves
this restriction. In each of the other replacements, every vertex
of $V(C_r)\setminus V(C')$ has more than $m_r$ neighbours in $J$.
It therefore suffices to check that these vertices lie in the
prescribed set.

\section{Two maximum independent sets}
\label{sec:maximum-independent}

Suppose that $G$ is a
non-Hamiltonian $3/2$-tough $2K_2$-free graph of order $n\geq 3$, and
set $\alpha=\alpha(G)$. 
Since $G$ is not complete, $\alpha\geq 2$.
Moreover, $n\ge 5$.
Otherwise if $n\in\{3,4\}$ and $x,y$ are nonadjacent, deleting
the other $n-2$ vertices contradicts $3/2$-toughness.

\subsection{Cycles of order \texorpdfstring{$n-1$}{n-1}}
\label{subsec:near-spanning-cycles}

We first 
show that every maximum independent set \(I\) is contained in a cycle of order \(n-1\).
\begin{lemma}\label{lem:n-1cycle}
For every maximum independent set $I$ of $G$, there is a cycle of
order $n-1$ containing $I$.
\end{lemma}

\begin{proof}
Let $I$ be a maximum independent set and   set $S=V(G)\setminus I$.
Since $G-S$ has $\alpha$ components, toughness gives
$|S|\geq 3\alpha/2>\alpha$.
Complete $S$ to a clique, obtaining a split graph $G^+$.
For every $X\subseteq V(G)$, we have
$\omega(G^+-X)\leq\omega(G-X)$, so $G^+$ remains $3/2$-tough.
By Theorem \ref{thm:KratschLehelMuller1996}, $G^+$ contains
a Hamilton cycle.

This cycle has $2\alpha$ edges between $I$ and $S$, and
$|S|-\alpha>0$ edges within $S$.
Delete all edges of the latter kind, including those belonging to $G$.
The remaining graph is a spanning disjoint union of paths in $G$.
Every vertex of $I$ still has degree two, so all path ends, including
singleton paths, belong to $S$.
Lemma~\ref{lem:path}\textup{(iii)}, applied with $J=I$, covers
$V(G)$ by at most $\alpha-|I|+1=1$ such path. Hence $G$ has a
Hamilton path with both endvertices outside $I$.

Write this path as
$P=v_0v_1\cdots v_{n-2}v_{n-1}$.
Since $G$ is non-Hamiltonian,
$v_0v_{n-1}\notin E(G)$.
The edges $v_0v_1$ and $v_{n-2}v_{n-1}$ are vertex-disjoint
because $n\geq5$.
Thus at least one of
$v_0v_{n-2}$, $v_1v_{n-1}$, and $v_1v_{n-2}$ is an edge.
If $v_0v_{n-2}\in E(G)$, then
$v_0Pv_{n-2}v_0$ is a cycle of order $n-1$.
If $v_1v_{n-1}\in E(G)$, then
$v_1Pv_{n-1}v_1$ is a cycle of order $n-1$.
Otherwise, $v_1v_{n-2}\in E(G)$, and
$v_1Pv_{n-2}v_1$ is a cycle of order $n-2\geq3$
with exactly $v_0$ and $v_{n-1}$ outside it. These vertices are
nonadjacent. Thus in every case there is a dominating cycle $Q$
containing $I$.
If $\ell$ is the maximum order of such a cycle,
Lemma~\ref{lem:cycle} gives
$\alpha\geq |I|+n-\ell-1$, and hence $\ell\geq n-1$.
Since $G$ is non-Hamiltonian, $\ell=n-1$.

\end{proof}


\subsection{Two maximum independent sets}
\label{subsec:extremal-pair}

Set
 $D=\{z\in V(G):G-z\text{ has a Hamilton cycle}\}$.
By Lemma~\ref{lem:n-1cycle}, $D\setminus I\neq\varnothing$
for every maximum independent set $I$.

Fix a maximum independent set $I$ and choose $u\in D\setminus I$ such
that
 $m=|N_I(u)|=\max_{z\in D\setminus I}|N_I(z)|$.
We have $m\geq 1$ by the maximality of $I$.
Since $I$ is independent, the choice of $u$    gives
\begin{equation}\label{eq:s3-D-bound}
 |N_I(z)|\leq m\qquad(z\in D).
\end{equation}


\begin{lemma}\label{prop:s3-omitted-vertices}
Let $Q$ be a dominating cycle of order $n-2$ containing $I$, and write
$V(G)\setminus V(Q)=\{r,s\}$.
If $|N_I(r)|>m$, then $s\in D$.
\end{lemma}

\begin{proof}
By \eqref{eq:s3-D-bound}, $r\notin D$.
Suppose that $s\notin D$ as well.
Let $\mathcal H$ consist of the dominating cycles $C$ containing $I$
for which $V(G)\setminus V(C)$ is disjoint from $D$.
Then $Q\in\mathcal{H}$.
The maximum order of a cycle in $\mathcal{H}$ is $n-2$: the given
cycle attains this order, a Hamilton cycle is excluded, and the
unique vertex outside a cycle of order $n-1$ would belong to $D$.
Choose $C_0\in\mathcal{H}$ first of maximum order and then with
$m_0=|N_I(U_0)|$ maximum, where $U_0=V(G)\setminus V(C_0)$.
Since $Q$    has maximum order and $|N_I(r)|>m$, we have $m_0>m$.

Use the extremal construction of Lemma~\ref{lem:cycle} with
$\mathcal F_0=\mathcal H$ and $J=I$. We verify the required closure
conditions in the order in which they are needed. Suppose that the
induction has reached $C_i$ with \eqref{eq:cycle-order} established up to index $i$.
In particular, $|C_i|=n-2-i$ and $m_i\ge m_0>m$.
Write $U_i=V(G)\setminus V(C_i)$ and $M_i=N_I(U_i)$.
Since $I$ is a maximum independent set and $U_i\ne\emptyset$,
we have $M_i\ne\emptyset$. By Fact~\ref{fact:2K2-property}(ii), choose $w\in U_i$
with $N_I(w)=M_i$. Orient $C_i$ and set
\[
 A_i=(M_i)^-_{C_i},\qquad
 T_i=\{x\in A_i:N_{U_i}(x)=\emptyset,\;
                         N_I(x)\setminus M_i\ne\emptyset\}.
\]
The independence of $I$ gives $A_i\cap I=\emptyset$.
For every $x\in T_i$, the nonedge $xw$ and Fact~\ref{fact:2K2-property}\textup{(ii)} give
$N_I(x)\supsetneq M_i$. Hence
 $|N_I(x)|>m_i\ge m_0>m$,
and therefore $x\notin D$ by \eqref{eq:s3-D-bound}.

A comparison cycle with vertex set $V(C_i)\cup X$, where
$\emptyset\ne X\subseteq U_i$, has outside vertex set $U_i\setminus X$.
A comparison cycle with vertex set
$(V(C_i)\setminus\{x\})\cup\{w\}$, where $x\in T_i$, has outside
vertex set $(U_i\setminus\{w\})\cup\{x\}$.
Both cycles are dominating and contain $I$, and their outside
vertex sets are disjoint from $D$. Thus these cycles belong to
$\mathcal H$ without any
use of Lemma~\ref{lem:predecessor-structure}. The first two extremal comparisons in the proof of
Lemma~\ref{lem:cycle} consequently exclude them. The cycle case of Lemma~\ref{lem:predecessor-structure} now applies to $C_i$,
with $J=I$, $u=w$, and its set $D$ taken to be $A_i$.
In particular, $A_i$ is independent.

The remaining two constructions have outside vertex sets
\[
 (U_i\setminus\{w,y\})\cup\{x\}
 \quad\text{and}\quad
 (U_i\setminus\{w\})\cup\{x,z\},
\]
respectively, where $x,z\in T_i$ are distinct in the second
construction and $y\in U_i\setminus\{w\}$ in the first.
Both constructed cycles are dominating, since $C_i$ is dominating,
$T_i\subseteq A_i$ is independent, and $T_i$ is anticomplete to $U_i$.
Their outside vertex sets are disjoint from $D$, since
$T_i\cap D=\emptyset$. The constructed cycles also contain $I$,
so they belong to $\mathcal H$.
The remaining steps of Lemma~\ref{lem:cycle} therefore either give the required
independent-set bound or establish \eqref{eq:cycle-order} at index $i+1$.
This verifies the closure hypothesis of Corollary~\ref{claim:cycle-restricted} throughout
the induction. Applying the corollary with $\ell_{\mathcal H}=n-2$
gives
 $\alpha\ge |I|+n-(n-2)-1=\alpha+1$,
a contradiction.
\end{proof}

Let $C$ be any oriented Hamilton cycle of $G-u$, and set
$I_0=I\setminus N_I(u)$. Since $I$ is independent,
$N_I(u)_C^-\subseteq N_G(u)_C^-\setminus I$ and
$|N_I(u)_C^-|=m$.

\begin{lemma}\label{lem:s3-structure}
The following assertions hold.
\begin{enumerate}[label=\textup{(\roman*)},nosep]
\item[\textup{(i)}]
The set $N_G(u)_C^+\cup\{u\}$ is independent. In particular,
$d_G(u)\leq\alpha-1$ (see \cite[Lemma~3.3 and its proof]{OtaSanka2022}).

\item[\textup{(ii)}]
Exactly one vertex $a\in N_I(u)_C^-$ has a neighbour in $I_0$.
Moreover, $N_I(a)\supsetneq N_I(u)$,
$N_G(u)_C^-\setminus I=N_I(u)_C^-$,
$N_I(u)_C^-\setminus D=\{a\}$, and
$ua_C^-\notin E(G)$.
Exactly one vertex of $N_I(u)_C^+$ has a neighbour in $I_0$ as well.

\item[\textup{(iii)}]
For every $x\in N_G(u)\setminus I$, both $x_C^-$ and $x_C^+$ belong
to $I_0$.
\end{enumerate}
\end{lemma}

\begin{proof}
All predecessors and successors in the proof are taken on $C$.
By \cite[Proposition~1.6]{OtaSanka2022}, $G$ has a $2$-factor.
Choose a $2$-factor $F$ of $G$ with the minimum number of components.
Since $G$ is non-Hamiltonian, $\omega(F)\geq2$, whereas $F'=C$
is a $2$-factor of $G-u$ with $\omega(F')=1$.
Thus $u$ is co-absorbable with respect to $F$.
Taking $F'=C$ in \cite[Lemma~3.3 and its proof]{OtaSanka2022}
gives \textup{(i)}, with successors taken on $C$.
Applying the same result to the reverse orientation of $C$ shows
that $N_G(u)_C^-\cup\{u\}$ is independent as well. In particular,
$N_G(u)_C^-\cap N_G(u)=\varnothing$.

If no vertex of $N_I(u)_C^-$ has a neighbour in $I_0$, then
$\{u\}\cup N_I(u)_C^-\cup I_0$ is independent of order
$\alpha+1$.
Choose a vertex $a\in N_I(u)_C^-$ with a neighbour $x\in I_0$.
For every $y\in N_I(u)$, the edges $ax,uy$ have no cross edges
$au,xu,xy$.
Thus $ay\in E(G)$, proving $N_I(a)\supsetneq N_I(u)$.
It follows from \eqref{eq:s3-D-bound} that $a\notin D$.

We next establish the predecessor relations needed to determine the
remaining vertices of $N_G(u)_C^-\setminus I$.
For every $x\in N_G(u)_C^-\setminus D$, we have
\begin{equation}\label{eq:s3-predecessor-switch}
    ux^-\notin E(G),\qquad
    N_{N_G(u)_C^-}(x^-)=\{x\}.
\end{equation}
Indeed,   set $y=x^+$.
If $ux^-\in E(G)$, replacing $x^-xy$ on $C$ by $x^-uy$ gives
a Hamilton cycle of $G-x$, contrary to $x\notin D$.
Hence $x^-\notin N_G(u)$.
Also, $x^-\notin N_G(u)_C^-$, since $x^-x$ is an edge and
$N_G(u)_C^-$ is independent.

Suppose that $x^-z\in E(G)$ for some
$z\in N_G(u)_C^-\setminus\{x\}$, and   set $t=z^+$.
The cycle $x^-z\overleftarrow{C}yutCx^-$ has vertex set
$V(G)\setminus\{x\}$.
To see that it is simple, note that the order on $C$ is
$x^-,x,y,\ldots,z,t,\ldots,x^-$.
The vertex $x^-$ lies outside
$N_G(u)\cup N_G(u)_C^-$, whereas
$x,z\in N_G(u)_C^-$ and $y,t\in N_G(u)$, so these five vertices are distinct; $u$ lies outside $C$.
The arcs $yCz$ and $tCx^-$ are disjoint and cover
$V(C)\setminus\{x\}$.
This   contradicts $x\notin D$, proving
\eqref{eq:s3-predecessor-switch}.

For distinct $x,z\in N_G(u)_C^-\setminus D$, the vertices
$x,x^-,z,z^-$ are distinct, and the edges $xx^-,zz^-$ have no
cross edges $xz,xz^-,zx^-$.
Consequently, $x^-z^-\in E(G)$ by $2K_2$-freeness.

Suppose now that
$x\in\bigl(N_G(u)_C^-\setminus I\bigr)\setminus\{a\}$ and
$x\notin D$.
Since $a,x\in N_G(u)_C^-\setminus D$, the preceding paragraph gives
$a^-x^-\in E(G)$.
Writing $y=a^+$ and $z=x^+$, we obtain the cycle
$a^-x^-\overleftarrow{C}yuzCa^-$.
The six old vertices $a^-,a,y,x^-,x,z$ are distinct: the vertices
$a^-,x^-$ lie outside $N_G(u)\cup N_G(u)_C^-$, while
$a,x\in N_G(u)_C^-$ and $y,z\in N_G(u)$.
Its two old arcs $yCx^-$ and $zCa^-$ are disjoint and partition
$V(C)\setminus\{a,x\}$.
It is therefore a simple cycle of order $n-2$ containing $I$, with
$V(G)\setminus V(Q)=\{a,x\}$, where $Q$ denotes this cycle.
As $ax\notin E(G)$, $Q$ is dominating. Since $|N_I(a)|>m$,
Lemma~\ref{prop:s3-omitted-vertices} gives $x\in D$, a
contradiction.
Thus
$\bigl(N_G(u)_C^-\setminus I\bigr)\setminus\{a\}\subseteq D$.

If a vertex $x$ in this set had a neighbour $y\in I_0$, then, for
every $z\in N_I(u)$, the edges $xy,uz$ and the nonedges
$xu,yu,yz$ would force $xz\in E(G)$.
This would give $|N_I(x)|>m$, contrary to
\eqref{eq:s3-D-bound}.
Hence
$\bigl(N_G(u)_C^-\setminus I\bigr)\setminus\{a\}$
is anticomplete to $I_0$, and
$I_0\cup\bigl((N_G(u)_C^-\setminus I)\setminus\{a\}\bigr)
\cup\{u\}$ is independent.
Its order is $\alpha-m+|N_G(u)_C^-\setminus I|$, so
$|N_G(u)_C^-\setminus I|\leq m$.
Together with
$N_I(u)_C^-\subseteq N_G(u)_C^-\setminus I$ and
$|N_I(u)_C^-|=m$, this yields
$N_G(u)_C^-\setminus I=N_I(u)_C^-$ and
$N_I(u)_C^-\setminus D=\{a\}$.
It    shows that $a$ is the unique vertex of $N_I(u)_C^-$ with
a neighbour in $I_0$.
The nonedge $ua^-\notin E(G)$ follows from
\eqref{eq:s3-predecessor-switch}.

Finally, let $x\in N_G(u)\setminus I$.
If $x^-\notin I$, then
$x^-\in N_G(u)_C^-\setminus I=N_I(u)_C^-$, and hence
$x\in N_I(u)$, a contradiction.
Thus $x^-\in I$.
Since $x^-\in N_G(u)_C^-$ is nonadjacent to $u$, we have
$x^-\in I_0$.
The arguments above hold for either orientation of $C$.
Applying them to the reverse orientation shows that exactly one
vertex of $N_I(u)_C^+$ has a neighbour in $I_0$ and that
$x_C^+\in I_0$.
This completes the proof of \textup{(ii)} and \textup{(iii)}.
\end{proof}

Among all pairs $(I,u)$ with $I$ a maximum independent set and
$u\in D\setminus I$, choose one for which $|I\cap D|$ is maximum and,
subject to this, $|N_I(u)|$ is maximum.
Lemma~\ref{lem:n-1cycle} guarantees that such pairs exist.
This choice    attains the fixed-$I$ maximum used in
\eqref{eq:s3-D-bound}.
Choose any oriented Hamilton cycle $C$ of $G-u$, and retain the
notation $m,I_0,a$ above.
  Set $K_0=\{u\}\cup\bigl(N_I(u)_C^-\setminus\{a\}\bigr)$ and
$I'=I_0\cup K_0$.

\begin{lemma}\label{lem:secindependent}
The set $I'$ is a maximum independent set.
Moreover, $N_I(u)\subseteq D$, and there is a vertex
$v\in N_I(u)$ such that $N_{I'}(v)=K_0$ and $(I',v)$ attains the
same two maxima as $(I,u)$.
\end{lemma}

\begin{proof}
By Lemma~\ref{lem:s3-structure}, the set $K_0$ is independent,
anticomplete to $I_0$, and has order $m$.
Also, $K_0\subseteq D$, since $u\in D$ and
$N_I(u)_C^-\setminus\{a\}\subseteq D$.
Thus $I'$ is independent of order $|I_0|+m=\alpha$.
Its intersection with $D$ is obtained from $I\cap D$ by deleting
$N_I(u)\cap D$ and adding the $m$ vertices of $K_0$.
Since Lemma~\ref{lem:n-1cycle}    gives $D\setminus I'\neq\varnothing$,
the first maximum in the choice of $(I,u)$ forces
$N_I(u)\subseteq D$ and $|I'\cap D|=|I\cap D|$.

The sets $N_{K_0}(z)$, for $z\in N_I(u)$, are linearly ordered by
inclusion.
Otherwise, two vertices of $N_I(u)$ with incomparable
neighbourhoods, together with one vertex in each set difference,
would induce a $2K_2$, since both $N_I(u)$ and $K_0$ are independent.
Their union is $K_0$: the vertex $u$ is adjacent to every vertex of
$N_I(u)$, and every vertex of
$N_I(u)_C^-\setminus\{a\}$ has its cycle successor in $N_I(u)$.
Hence some $v\in N_I(u)$ is adjacent to all of $K_0$.
Since $v\in I$, it has no neighbour in $I_0$, and therefore
$N_{I'}(v)=K_0$.

As $N_I(u)\subseteq D$ and $N_I(u)\cap I'=\varnothing$, we have
$v\in D\setminus I'$.
Moreover, $|I'\cap D|=|I\cap D|$, so the second choice of $(I,u)$
gives $|N_{I'}(z)|\leq m$ for every $z\in D\setminus I'$.
Equality holds for $v$, proving that $(I',v)$ attains the same two
maxima as $(I,u)$.
\end{proof}


Fix such a vertex $v$.
In particular, $uv\in E(G)$.
Lemma~\ref{prop:s3-omitted-vertices} and
Lemma~\ref{lem:s3-structure} apply to the pair $(I',v)$, with the
same value of $m$.
The sets corresponding to $N_I(u)$ and $I_0$ are $K_0$ and
$I'\setminus K_0=I_0$, respectively; in
Lemma~\ref{lem:s3-structure}, the cycle may be any oriented Hamilton
cycle $C'$ of $G-v$.
Set $k=|I_0|$, so $\alpha=m+k$, and   set
$T=N_G(u)\setminus N_I(u)$ and $T'=N_G(v)\setminus K_0$.
We have $k\geq 1$, since $a$ has a neighbour in $I_0$.
Applying Lemma~\ref{lem:s3-structure}\textup{(i)} to both pairs gives
$|T|\leq k-1$ and $|T'|\leq k-1$.

\section{Proof of Theorem~\ref{thm:main}}
\label{sec:proof}

\begin{proof}[Proof of Theorem~\ref{thm:main}]
Suppose, to the contrary, that $G$ is a non-Hamiltonian graph satisfying
the hypotheses.
Choose the pair $(I,u)$ by maximizing $|I\cap D|$ over all pairs
with $I$ a maximum independent set and $u\in D\setminus I$, and
then maximizing $|N_I(u)|$ over all pairs attaining the first maximum.
Take the oriented Hamilton cycle $C$ of $G-u$ and the sets
$I_0,K_0,I'$ used in the construction preceding
Lemma~\ref{lem:secindependent}, and choose $v$ as in that lemma.
In particular, $C$ is fixed together with these sets and
\[
 I_0=I\cap I',\qquad
 K_0=\{u\}\cup\bigl(N_I(u)_C^-\setminus\{a\}\bigr),\qquad
 I'=I_0\cup K_0,
\]
where $a$ is the unique vertex of $N_I(u)_C^-$ with a neighbour
in $I_0$. We have $N_{I'}(v)=K_0$, $|N_I(u)|=|K_0|=m$,
$uv\in E(G)$, and $(I',v)$ attains the same two maxima.
Choose any oriented Hamilton cycle $C'$ of $G-v$.

Retain $T=N_G(u)\setminus N_I(u)$ and
$T'=N_G(v)\setminus K_0$, and   set
$S=N_G(I_0)$ and
$p=|N_I(u)_C^-\cap N_I(u)_C^+|$.

  Set $x=a_C^-$ and $y=a_C^+\in N_I(u)$.
By Lemma~\ref{lem:s3-structure}\textup{(ii)}, $a\notin D$,
$N_I(a)\supsetneq N_I(u)$, and $ux\notin E(G)$; in particular,
$x\notin N_I(u)$.
Deleting $a$ from $C$ and adding $yu$ gives the Hamilton path
$P_0=x\overleftarrow{C}yu$ of $G-a$.

We rotate paths obtained from $P_0$ while keeping $x$ fixed.
For a path $P=v_0v_1\cdots v_s$ with $v_0=x$ and $v_s=w$,
and an edge $wv_i$ with $1\leq i\leq s-2$, replace
$v_iv_{i+1}$ by $v_iw$ to obtain the path
$v_0v_1\cdots v_iwv_{s-1}\cdots v_{i+1}$.
Let $W$ be the set of endvertices other than $x$ of all paths obtained
from $P_0$ by a sequence of these rotations, including $u$.

Let $P$ be any such path, with ends $x,w$.
An edge $wa$ would give the Hamilton cycle $axPwa$ of $G$.
An edge $wx$ would close $P$ to a Hamilton cycle of $G-a$,
contrary to $a\notin D$.
Thus $W$ has no edges to $\{a,x\}$.
An edge within $W$, together with $ax$, would induce a $2K_2$;
hence $W$ is independent.
Consequently, each vertex of $W$ is isolated in $G-N_G(W)$,
and the edge $ax$ remains.
Let $H$ be the component of $G-N_G(W)$ containing $ax$.
Since $G$ is $2K_2$-free, every other component is an isolated vertex.
Let $h$ be the number of vertices of $P_0$ with two neighbours in $W$
along $P_0$.

\begin{claim}\label{claim:n}
    $n=2|W|-1-h+\alpha+|V(H)\setminus I|$.
\end{claim}

\begin{proof}
We first establish
\begin{equation}\label{eq:s4-rotation-properties}
 N_G(W)=N_{P_0}(W),\qquad
 W\setminus I=K_0,\qquad
 N_G(W)\cap I=N_I(u).
\end{equation}
The inclusion $N_{P_0}(W)\subseteq N_G(W)$ is immediate.
Suppose that $z\in N_G(W)\setminus N_{P_0}(W)$.
Then $z\notin W\cup\{a,x\}$.
If a rotation first changes an edge incident with $z$, then $z$ must
be its pivot, since both the old and the new endvertex belong to $W$.
The edge deleted at $z$ is still an edge of $P_0$, and its other end
becomes a member of $W$, contrary to $z\notin N_{P_0}(W)$.
Thus all path edges incident with $z$ remain unchanged in every
rotation sequence.
Choose a path $P$ obtained from $P_0$ whose end $w$ is adjacent to $z$.
If $wz$ is an edge of $P$, it is an edge of $P_0$, a contradiction.
Otherwise, a rotation at $z$ makes an original $P_0$-neighbour of $z$
an endvertex,   a contradiction.
This proves $N_G(W)=N_{P_0}(W)$.

For each $z\in N_G(u)\setminus\{y\}$, a rotation of $P_0$ at $z$ has
end $z_C^-$.
These rotations are possible because $ux\notin E(G)$ and $y$ is the
only neighbour of $u$ along $P_0$.
Since $u\in W$ and
$K_0=\{u\}\cup(N_I(u)_C^-\setminus\{a\})$, we obtain
\begin{equation}\label{eq:s4-initial-rotations}
 K_0\cup T_C^-\subseteq W.
\end{equation}

Suppose that $w\in W\setminus I$ has a neighbour $z\in I_0$.
For every $t\in N_I(u)$, the edges $wz,ut$ have no cross edges
$wu,zu,zt$.
Hence $wt\in E(G)$, and $N_I(w)\supsetneq N_I(u)$.
By the extremal choice of $u$, we have $w\notin D$.
Take a path $P$ obtained from $P_0$ with ends $x,w$, and let $tw$
be its last edge.
The edges $ax,tw$ have four distinct ends, since
$|V(P)|=n-1\geq4$.
As $aw,xw\notin E(G)$, at least one of $at,xt$ is an edge.
If $at\in E(G)$, then $axPt a$ is a Hamilton cycle of $G-w$,
contrary to $w\notin D$.
Otherwise, $xPtx$ is a cycle of order $n-2$ containing $I$.
The two vertices outside this cycle are $a,w$; they are nonadjacent,
so the cycle is dominating. Also, $|N_I(a)|>m$.
Lemma~\ref{prop:s3-omitted-vertices} therefore gives $w\in D$,
  a contradiction.

It follows that $W\setminus I$ is anticomplete to $I_0$.
Also, $W\cap N_I(u)=\varnothing$, since $u\in W$ is adjacent to
every vertex of $N_I(u)$.
Thus $I_0\cup(W\setminus I)$ is independent, so
$|W\setminus I|\leq\alpha-k=m$.
Together with \eqref{eq:s4-initial-rotations} and $|K_0|=m$, this gives
$W\setminus I=K_0$.
No vertex of $W$ has a neighbour in $I_0$, whereas $u$ is adjacent to
every vertex of $N_I(u)$.
Hence $N_G(W)\cap I=N_I(u)$, completing the proof of
\eqref{eq:s4-rotation-properties}.
In particular, $W\subseteq I'$.

The independent set $W$ contains exactly one end of $P_0$.
Counting the edges of $P_0$ incident with $W$ gives $2|W|-1$ edges.
Each vertex of $N_{P_0}(W)$ is counted once, except that the $h$
vertices with two neighbours in $W$ are counted twice.
By \eqref{eq:s4-rotation-properties},
$|N_G(W)|=2|W|-1-h$.

Every vertex of $I'=I_0\cup K_0$ remains in $G-N_G(W)$.
An isolated vertex outside $I'$ could be added to this maximum
independent set, so every isolated vertex belongs to $I'$.
Moreover, $N_I(u)\subseteq N_G(W)$ and every vertex of
$K_0\subseteq W$ is isolated in $G-N_G(W)$.
Consequently, $I'\cap V(H)=I\cap V(H)$, and the vertices outside
$N_G(W)\cup I'$ are precisely $V(H)\setminus I$.
Therefore
\[
 n=|N_G(W)|+\alpha+|V(H)\setminus I|
  =2|W|-1-h+\alpha+|V(H)\setminus I|,
\]
as required.
\end{proof}

\begin{claim}\label{claim:h}
    $h\geq p+|T|+|W|-\alpha$.
\end{claim}

\begin{proof}
There are $p$ oriented subpaths $v_0v_1v_2$ of $C$ with
$v_0,v_2\in N_I(u)$.
For each such subpath,
$v_1\in N_I(u)_C^-\cap N_I(u)_C^+$ and $v_1\neq a$, since
$a_C^-=x\notin N_I(u)$.
Thus $v_1\in K_0\subseteq W$.
If $v_0\neq y$, then
$(v_0)_C^-\in N_I(u)_C^-\setminus\{a\}\subseteq W$,
and both cycle edges incident with $v_0$ remain on $P_0$.
If $v_0=y$, its neighbour $a$ on $C$ is replaced on $P_0$ by
$u\in W$.
In either case, $v_0$ has two neighbours in $W$ along $P_0$.
Distinct subpaths give distinct vertices $v_0$, accounting for $p$
vertices of $N_I(u)$ counted by $h$.

By Lemma~\ref{lem:s3-structure}\textup{(iii)} and
\eqref{eq:s4-initial-rotations},
$T_C^-\subseteq W\cap I_0$ and $T_C^+\subseteq I_0$.
If $t\in T$ and $t_C^+\in W$, then both neighbours of $t$ on $C$
belong to $W$.
Its incident cycle edges remain on $P_0$: the vertex $t$ differs from
$a,x$, which are nonadjacent to $u$, and from $y\in I$.
Thus $t$ is    counted by $h$.
These vertices lie outside $I$ and are distinct from those already
counted in $N_I(u)$.
Their number is at least
 $|T_C^+\cap W|
 \geq |T|+|W\cap I_0|-k
 =|T|+|W|-\alpha,$
where the equality follows from
$W\cap N_I(u)=\varnothing$, $W\setminus I=K_0$, and $|K_0|=m$.
Adding the two counts proves the claim.
\end{proof}

\begin{claim}\label{claim:|W|+|V(H)setminus I|}
    $|W|+|V(H)\setminus I|\leq\alpha-1$.
\end{claim}

\begin{proof}
Let $z\in V(H)\setminus I$.
As shown in Claim~\ref{claim:n},
$I'\cap V(H)=I\cap V(H)$, so $z\notin I'$.
The maximality of $I'$ gives a neighbour $y\in I'$.
Since all of $I'$ remains in $G-N_G(W)$, we have
$y\in I'\cap V(H)=I\cap V(H)\subseteq I_0$.
For any $t\in N_I(u)$, the edges $zy,ut$ have no cross edges
$zu,yu,yt$: the first is absent because $u\in W$ and
$z\notin N_G(W)$, the second because $y\in I_0$, and the third
because $I$ is independent.
Hence $zt\in E(G)$.
In particular, $z$ is adjacent to $v\in N_I(u)$.
As $K_0\cap V(H)=\varnothing$, it follows that $z\in T'$.

By Lemma~\ref{lem:s3-structure}\textup{(iii)} applied to
$(I',v,C')$, both neighbours of $z$ on $C'$ belong to $I_0$.
They    lie in $H$, since all vertices of $I_0$ survive the deletion
of $N_G(W)$.
Thus $C'[V(H)]$ is bipartite with parts $I\cap V(H)$ and
$V(H)\setminus I$, and every vertex of the latter part has degree two.
Since $v\in N_I(u)\subseteq N_G(W)$, all of $H$ lies on $C'$.
The vertex $u$ lies on $C'$ but outside $H$, so $C'[V(H)]$ is a
nonempty forest with exactly $2|V(H)\setminus I|$ edges.
Consequently,
 $|I\cap V(H)|-|V(H)\setminus I|
 =\omega\bigl(C'[V(H)]\bigr)\geq1$.
The disjoint sets $W$ and $I\cap V(H)$ are both contained in $I'$.
Therefore
$ |W|+|V(H)\setminus I|
 \leq |W|+|I\cap V(H)|-1
 \leq\alpha-1,$
proving the claim.
\end{proof}

\medskip
Combining Claims~\ref{claim:n}--\ref{claim:|W|+|V(H)setminus I|},
we obtain
\[
\begin{aligned}
 n
 &=2|W|-1-h+\alpha+|V(H)\setminus I|\\
 &\leq |W|+2\alpha-1-p-|T|+|V(H)\setminus I|\\
 &\leq 3\alpha-2-p-|T|.
\end{aligned}
\]
By Lemma~\ref{lem:s3-structure}\textup{(ii)}, each of
$N_I(u)_C^-$ and $N_I(u)_C^+$ contains exactly one vertex with
a neighbour in $I_0$.
Their union has $2m-p$ vertices and is contained in
$V(C)\setminus I$.
Also, $S\subseteq V(C)\setminus I$, since $I$ is independent and
$u$ has no neighbour in $I_0$.
Thus
\[
\begin{aligned}
 |S|
 &\leq |V(C)\setminus I|-(2m-p)+2\\
 &=n-\alpha-2m+p+1\\
 &\leq 2\alpha-2m-1-|T|\\
 &=2k-1-|T|.
\end{aligned}
\]

We next show that $S=T\cup T'$.
By Lemma~\ref{lem:s3-structure}\textup{(iii)}, every vertex of $T$
has both neighbours on $C$ in $I_0$.
Applying the same result to $(I',v,C')$ shows that every vertex of
$T'$ has both neighbours on $C'$ in $I_0$.
Hence $T\cup T'\subseteq S$.
Conversely, let $y\in I_0$ and $z\in N_G(y)$.
Since $yu,yv\notin E(G)$, the edges $yz,uv$ force
$zu\in E(G)$ or $zv\in E(G)$.
Neither $N_I(u)$ nor $K_0$ has an edge to $I_0$, so
$z\in T\cup T'$.
This proves $S=T\cup T'$.

The degree bound in Lemma~\ref{lem:s3-structure}\textup{(i)},
applied to $(I',v)$, gives
$|T'|=d_G(v)-m\leq\alpha-1-m=k-1$.
Therefore
$ |S|\leq |T|+|T'|\leq |T|+k-1$.
Adding the two bounds on $|S|$ yields $2|S|\leq3k-2$.
On the other hand, every vertex of $I_0$ is isolated in $G-S$,
while the edge $uv$ remains.
Hence $\omega(G-S)\geq k+1\geq2$, and the toughness condition gives
$
 2|S|\geq3\omega(G-S)\geq3(k+1)$,
a contradiction.
\end{proof}





\begin{thebibliography}{99}
\bibitem{BauerBroersmaSchmeichel2006}
D. Bauer, H. J. Broersma, and E. Schmeichel, Toughness in graphs---A survey, \emph{Graphs Combin.} \textbf{22} (2006), no. 1, 1--35.

\bibitem{BauerBroersmaVeldman2000}
D. Bauer, H. J. Broersma, and H. J. Veldman, Not every $2$-tough graph is Hamiltonian, \emph{Discrete Applied Mathematics} \textbf{99} (2000), no.~1--3, 317--321.



\bibitem{Chvatal1973}
V. Chv{\'a}tal, Tough graphs and hamiltonian circuits, \emph{Discrete Mathematics} \textbf{5} (1973), no.~3, 215--228.


\bibitem{KratschLehelMuller1996}
Dieter Kratsch, Jen{\H o} Lehel, and Haiko M{\"u}ller, Toughness, hamiltonicity and split graphs, \emph{Discrete Mathematics} \textbf{150} (1996), no.~1--3, 231--245.

\bibitem{BroersmaPatelPyatkin2014}
Hajo Broersma, Viresh Patel, and Artem Pyatkin, On toughness and hamiltonicity of {$2K_2$}-free graphs, \emph{Journal of Graph Theory} \textbf{75} (2014), no.~3, 244--255.

\bibitem{Shan2020}
Songling Shan, {Hamiltonian} cycles in {$3$}-tough {$2K_2$}-free graphs, \emph{Journal of Graph Theory} \textbf{94} (2020), no.~3, 349--363.

\bibitem{OtaSanka2022}
Katsuhiro Ota and Masahiro Sanka, {Hamiltonian} cycles in {$2$}-tough {$2K_2$}-free graphs, \emph{Journal of Graph Theory} \textbf{101} (2022), no.~4, 769--781.



\end{thebibliography}

\end{document}
